\documentclass[a4paper]{article}
\usepackage{graphicx}
\usepackage{amsmath,amssymb}
\usepackage{hyperref}
\usepackage{minted}
\usepackage{multirow}
\usepackage{float}
\usepackage{todonotes}
\usepackage{forest}
\usepackage{xstring}
\input{/home/donkarlo/Dropbox/repo/nd_language_ai_project/src/nd_language_ai/formal/computer/programming/tex/latex/code_clips/external_dot.tex}
\title{}
\date{}

\begin{document}
    \maketitle
    
    
    \section{Punkte und Teilergebnisse}
        
        \textbf{Deutsch:} Entscheidend für den Prüfungserfolg sind die Teilergebnisse folgender Subtests.
        \textbf{English:} The partial results of the following subtests are decisive for success in the exam.
        
        \begin{enumerate}
            \item Hören/Lesen \quad (Listening/Reading)
            \item Schreiben \quad (Writing)
            \item Sprechen \quad (Speaking)
        \end{enumerate}
        
        \subsection{Ermittlung des Teilergebnisses „Hören/Lesen“}
            
            \textbf{Deutsch:} Die Subtests „Hören“ und „Lesen“ bestehen aus insgesamt 45 Aufgaben. Für jede richtig gelöste Aufgabe erhalten die Teilnehmer/innen einen Punkt, so dass eine maximale Punktzahl von 45 erreicht werden kann.
            
            \textbf{English:} The subtests “Listening” and “Reading” consist of a total of 45 tasks. For each correctly solved task, participants receive one point, so the maximum possible score is 45 points.
            
            \textbf{Deutsch:} Für das Erreichen der Stufen A2 und B1 gilt:
            \textbf{English:} The following score ranges apply for reaching levels A2 and B1:
            
            \begin{table}[H]
                \centering
                \begin{tabular}{|c|c|}
                    \hline
                    Punkte (Points) & Stufen nach GER (CEFR Level) \\
                    \hline
                    33--45          & B1                           \\
                    20--32          & A2                           \\
                    0--19           & unter A2 (below A2)          \\
                    \hline
                \end{tabular}
            \end{table}
        
      
    % \bibliographystyle{plain}
    % \bibliography{/home/donkarlo/Dropbox/repo/shared_research/bibliography/bibliography.bib}
\end{document}